\documentclass[11pt,a4paper]{article}

\usepackage{amsmath,amssymb,amsthm}
\usepackage{mathtools}
\usepackage{bm}
\usepackage{enumerate}
\usepackage{hyperref}
\usepackage{cleveref}
\usepackage[margin=1in]{geometry}
\usepackage{booktabs}

% Theorem environments
\newtheorem{theorem}{Theorem}[section]
\newtheorem{lemma}[theorem]{Lemma}
\newtheorem{proposition}[theorem]{Proposition}
\newtheorem{corollary}[theorem]{Corollary}
\theoremstyle{definition}
\newtheorem{definition}[theorem]{Definition}
\newtheorem{example}[theorem]{Example}
\newtheorem{remark}[theorem]{Remark}

% Notation shortcuts
\newcommand{\E}{\mathbb{E}}
\newcommand{\Var}{\mathrm{Var}}
\newcommand{\Cov}{\mathrm{Cov}}
\newcommand{\Normal}{\mathcal{N}}
\newcommand{\IG}{\mathrm{IG}}
\newcommand{\given}{\,|\,}
\newcommand{\bpsi}{\boldsymbol{\psi}}
\newcommand{\btheta}{\boldsymbol{\theta}}
\newcommand{\bmu}{\boldsymbol{\mu}}
\DeclareMathOperator{\CE}{CE}
\DeclareMathOperator{\RDU}{RDU}

\title{Prior Predictive Specification for\\Bayesian Hierarchical Models:\\
Extending the Moment-Matching Framework to Observable Quantities}
\author{Technical Note}
\date{February 2026}

\begin{document}

\maketitle

\begin{abstract}
The current paper derives hyperpriors by matching the moments of the
unconditional distribution of parameters $\theta_i$ to a researcher-specified
``direct prior.'' This note develops an extension where the researcher instead
specifies beliefs about observable quantities---such as certainty
equivalents---and hyperpriors are derived so that the prior predictive
distribution matches these beliefs. When the mapping from $p$ preference
parameters to $p$ observable quantities is invertible, the moment-matching
framework applies directly in the observable space with no approximation:
the researcher specifies a direct prior over certainty equivalents, the
standard $\rho/\alpha_0$ decomposition derives hyperpriors in that space, and
posterior sampling proceeds in parameter space using only the forward mapping
and a Jacobian correction. We establish the conditions under which the
certainty equivalent mapping is globally injective---showing that prize
diversity across reference lotteries is the key structural requirement---and
demonstrate that the identity case (parameter-level specification) is a special
case with unit Jacobian.
\end{abstract}

\tableofcontents
\newpage

%==============================================================================
\section{Motivation}
%==============================================================================

The current paper asks the researcher to specify a direct prior over the
parameter vector $\btheta_i$ of a single individual:
\begin{equation}
    \theta_i^j \sim \Normal(\mu^{j*}, \sigma^{2j*}), \quad j = 1, \ldots, p
\end{equation}
The hyperprior is then derived so that the unconditional (marginal) distribution
of $\theta_i^j$ matches this direct prior. This works well when the researcher
has good intuition about the parameter space---for example, a CRRA coefficient
$r$ or the log of a Prelec curvature parameter.

However, researchers often have stronger intuition about \emph{observable
consequences} of preferences than about the parameters themselves. A researcher
might confidently say:
\begin{quote}
    ``For a lottery paying \$10, \$30, or \$50 with equal probability, I expect
    a typical person's certainty equivalent to be around \$25, and it would be
    unusual for it to fall below \$5 or above \$45.''
\end{quote}
This statement is about $g(\btheta)$---the certainty equivalent as a function
of risk preference parameters---not about $\btheta$ directly. The certainty
equivalent integrates information from all four parameters ($r$, $\phi$,
$\eta$, $\lambda$) through a nonlinear mapping, making it difficult to
``reverse engineer'' parameter-level beliefs from predictive-level beliefs.

This extension is motivated by three considerations:

\begin{enumerate}
    \item \textbf{Cognitive tractability.} Researchers may find it easier to
    reason about quantities in the same units as experimental outcomes (dollars,
    probabilities of choice) than about abstract model parameters.

    \item \textbf{Literature alignment.} The prior predictive checking
    literature (e.g., Box 1980, Gabry et al.\ 2019) emphasizes evaluating
    priors through their implications for observables. Specifying the prior
    \emph{directly} in predictive space makes this the starting point rather
    than a post-hoc check.

    \item \textbf{Cross-parameter coherence.} When $g$ depends on multiple
    parameters jointly, specifying beliefs about $g(\btheta)$ implicitly
    constrains the joint distribution of $\btheta$ in a way that
    parameter-by-parameter specification does not. This is relevant for the RDU
    model where the certainty equivalent depends on all four parameters
    simultaneously.
\end{enumerate}

%==============================================================================
\section{General Framework}
%==============================================================================

\subsection{Setup}

Consider a hierarchical model where:
\begin{itemize}
    \item $\btheta_i \in \mathbb{R}^p$ is the parameter vector for subject $i$
    \item $\btheta_i \given \bmu, \Sigma \sim \Normal(\bmu, \Sigma)$ (or with
        independent components: $\theta_i^j \given \mu_j, \sigma_j^2 \sim
        \Normal(\mu_j, \sigma_j^2)$)
    \item $(\bmu, \Sigma) \sim p(\bmu, \Sigma \given \bpsi)$ for some
        hyperprior with parameters $\bpsi$
    \item $g: \mathbb{R}^p \to \mathbb{R}$ is a known function mapping
        parameters to an observable quantity
\end{itemize}

\begin{definition}[Direct Prior Predictive]
The \emph{direct prior predictive} for $g$ is the unconditional distribution of
$Y_i = g(\btheta_i)$, integrating over all levels of the hierarchy:
\begin{equation}
    \label{eq:prior_predictive}
    p(Y_i) = \int p(Y_i \given \btheta_i) \, p(\btheta_i \given \bmu, \Sigma)
    \, p(\bmu, \Sigma \given \bpsi) \, d\btheta_i \, d\bmu \, d\Sigma
\end{equation}
When $Y_i = g(\btheta_i)$ is a deterministic function of $\btheta_i$, this
simplifies to:
\begin{equation}
    p(Y_i) = \int \delta(Y_i - g(\btheta_i)) \, p(\btheta_i \given \bpsi) \,
    d\btheta_i
\end{equation}
where $p(\btheta_i \given \bpsi) = \int p(\btheta_i \given \bmu, \Sigma) \,
p(\bmu, \Sigma \given \bpsi) \, d\bmu \, d\Sigma$ is the unconditional
(marginal) distribution of $\btheta_i$.
\end{definition}

The researcher specifies target moments for the prior predictive:
\begin{align}
    \E[g(\btheta_i)] &= m^* \\
    \Var[g(\btheta_i)] &= v^*
\end{align}
The task is to find hyperprior parameters $\bpsi$ such that the implied prior
predictive of $g$ matches these targets.

\subsection{Law of Total Variance Applied to \texorpdfstring{$g(\btheta)$}{g(theta)}}

The laws of iterated expectations and total variance apply to $g(\btheta_i)$
just as they do to $\theta_i$ itself. Define the conditional moments:
\begin{align}
    \label{eq:cond_mean}
    m(\bmu, \Sigma) &\coloneqq \E[g(\btheta_i) \given \bmu, \Sigma] \\
    \label{eq:cond_var}
    v(\bmu, \Sigma) &\coloneqq \Var[g(\btheta_i) \given \bmu, \Sigma]
\end{align}
These are the mean and variance of $g(\btheta_i)$ for a \emph{known} population
distribution $\Normal(\bmu, \Sigma)$, before integrating over hyperprior
uncertainty.

The unconditional moments then decompose as:
\begin{align}
    \label{eq:iterated_g}
    m^* = \E[g(\btheta_i)] &= \E_{\bpsi}[m(\bmu, \Sigma)] \\
    \label{eq:total_var_g}
    v^* = \Var[g(\btheta_i)] &= \underbrace{\E_{\bpsi}[v(\bmu, \Sigma)]}_{\text{average within-population variance}} + \underbrace{\Var_{\bpsi}[m(\bmu, \Sigma)]}_{\text{between-population variance}}
\end{align}
where $\E_{\bpsi}$ and $\Var_{\bpsi}$ denote expectations over the hyperprior.

\begin{remark}
Equations~\eqref{eq:iterated_g}--\eqref{eq:total_var_g} are completely general.
They hold for any hyperprior family, any function $g$, and any subject-level
distribution---they are the laws of iterated expectations and total variance
applied to $g(\btheta_i)$.
\end{remark}

The variance partition with $\rho$ generalizes directly:
\begin{align}
    \label{eq:rho_g}
    \rho \, v^* &= \Var_{\bpsi}[m(\bmu, \Sigma)] \\
    (1 - \rho) \, v^* &= \E_{\bpsi}[v(\bmu, \Sigma)]
\end{align}
The interpretation of $\rho$ is preserved: it is the fraction of prior
predictive variance attributable to uncertainty about the population center
(as measured through its effect on $g$), versus average within-population
dispersion.

\subsection{The Identity Case Recovers the Current Method}
\label{sec:identity}

When $g(\btheta_i) = \theta_i^j$ (the identity function applied to a single
parameter), the conditional moments simplify to:
\begin{align}
    m(\mu_j, \sigma_j^2) &= \E[\theta_i^j \given \mu_j, \sigma_j^2] = \mu_j \\
    v(\mu_j, \sigma_j^2) &= \Var[\theta_i^j \given \mu_j, \sigma_j^2] = \sigma_j^2
\end{align}
Substituting into \eqref{eq:total_var_g}:
\begin{equation}
    \sigma^{2j*} = \E[\sigma_j^2] + \Var[\mu_j]
\end{equation}
which is exactly the variance decomposition in the current paper. The $\rho$
decomposition yields the same moment targets:
$\Var[\mu_j] = \rho \sigma^{2j*}$ and $\E[\sigma_j^2] = (1-\rho)\sigma^{2j*}$.
The current method is the special case $g = \text{id}$.


%==============================================================================
\section{Why Invertibility Matters}
%==============================================================================

When $g: \mathbb{R}^p \to \mathbb{R}$ is a scalar-valued function of a vector
argument, the framework of Section~2 cannot be applied directly: a single
calibration quantity provides only 2 moment conditions ($m^*$, $v^*$) for $2p$
hyperprior unknowns. The natural resolution is to choose exactly $p$
calibration quantities so that
$\mathbf{g}: \mathbb{R}^p \to \mathbb{R}^p$ is an invertible mapping, and
work entirely in the calibration space.

With $p$ calibration quantities matching the $p$ preference parameters, we
have $2p$ moment conditions for $2p$ unknowns (after fixing $\rho$ and
$\alpha_0$)---exactly as in the identity case. Moreover, if $\mathbf{g}$ is
invertible, the prior can be specified directly in $\mathbf{g}$-space with no
approximation, and the nonlinearity of $\mathbf{g}$ is handled exactly by a
Jacobian correction during sampling. This is developed in the next section.


%==============================================================================
\section{The Invertible Case: Full Reparameterization}
\label{sec:invertible}
%==============================================================================

As noted in Section~3, the dimensionality problem is resolved by choosing $p$
calibration quantities to match the $p$ preference parameters, so that
$\mathbf{g}: \mathbb{R}^p \to \mathbb{R}^p$ is a \emph{diffeomorphism}
(smooth bijection with smooth inverse). With an invertible mapping, the
moment-matching framework applies directly in the predictive space.

\subsection{Change of Variables in the Prior}

Suppose we have $p$ preference parameters
$\btheta = (\theta^1, \ldots, \theta^p)$ and $p$ observable quantities
$\mathbf{Y} = (Y^1, \ldots, Y^p)$ defined by an invertible mapping:
\begin{equation}
    \mathbf{Y} = \mathbf{g}(\btheta), \qquad
    \btheta = \mathbf{g}^{-1}(\mathbf{Y})
\end{equation}
For the RDU application, $\btheta = (r, \phi, \eta)$ (the 3 preference
parameters, with noise $\lambda$ specified separately) and
$\mathbf{Y} = (\CE_1, \CE_2, \CE_3)$ are certainty equivalents for three
suitably chosen reference lotteries.

Instead of specifying a prior in parameter space, the researcher specifies a
prior directly in CE space:
\begin{equation}
    \label{eq:direct_prior_ce}
    Y_i^k \sim \Normal(m^{k*}, s^{2k*}), \qquad k = 1, \ldots, p
\end{equation}
This is the ``direct prior predictive''---it encodes beliefs about what the
observable preferences of a single random individual would look like. (The
Normal distribution is used here for continuity with the preceding development;
Section~\ref{sec:bounded} replaces it with a bounded distribution appropriate
for quantities like certainty equivalents.)

If we wish to evaluate the prior density in parameter space (e.g., because the
likelihood is naturally expressed in terms of $\btheta$), the standard change
of variables gives:
\begin{equation}
    \label{eq:jacobian_prior}
    p(\btheta) = p_{\mathbf{Y}}(\mathbf{g}(\btheta)) \;
    \left|\det J_{\mathbf{g}}(\btheta)\right|
\end{equation}
where $J_{\mathbf{g}}(\btheta)$ is the $p \times p$ Jacobian matrix with
entries $(J_{\mathbf{g}})_{kj} = \partial g_k / \partial \theta^j$, and
$p_{\mathbf{Y}}$ is the density specified in \eqref{eq:direct_prior_ce}. The
Jacobian corrects for the nonlinear warping of volume elements between the two
coordinate systems. This is the \emph{only} place the nonlinearity of
$\mathbf{g}$ appears.

\subsection{The BHM in Predictive Space}

The crucial insight is that the \emph{entire} Bayesian hierarchy---from
hyperprior through population distribution to individual-level prior---can be
specified in observable (CE) space. Model parameters $\btheta$ appear only in
the likelihood and through the Jacobian that links the two spaces. The
hyperprior, the population distribution, and the variance decomposition all
operate on certainty equivalents, denominated in the same units as experimental
outcomes.

Using Normal distributions for illustration (Section~\ref{sec:bounded}
develops the proper bounded treatment), the hierarchical model is:
\begin{align}
    \label{eq:bhm_ce_subject}
    Y_i^k \given \mu_k^Y, \sigma_k^{Y2} &\sim
        \Normal(\mu_k^Y, \sigma_k^{Y2}), \quad k = 1, \ldots, p \\
    \label{eq:bhm_ce_mean}
    \mu_k^Y &\sim \Normal(\mu_{0,k}^Y, \tau_{0,k}^{Y2}) \\
    \label{eq:bhm_ce_var}
    \sigma_k^{Y2} &\sim \IG(\alpha_0^Y, \beta_{0,k}^Y)
\end{align}
The moment-matching framework from the paper applies \emph{exactly} as before,
with no modifications:
\begin{align}
    m^{k*} &= \mu_{0,k}^Y \\
    s^{2k*} &= \E[\sigma_k^{Y2}] + \Var[\mu_k^Y]
        = \frac{\beta_{0,k}^Y}{\alpha_0^Y - 1} + \tau_{0,k}^{Y2}
\end{align}
The $\rho$ decomposition, the $\alpha_0$ specification, and the derived
hyperparameter formulas all carry over verbatim---just applied to the CE moments
$(m^{k*}, s^{2k*})$ rather than the parameter moments
$(\mu^{j*}, \sigma^{2j*})$.

\subsection{Posterior Sampling}
\label{sec:posterior_sampling}

The natural approach draws proposals for $\btheta_i$ directly in parameter
space. Each MCMC step requires:
\begin{enumerate}
    \item The \textbf{likelihood} $p(\text{data}_i \given \btheta_i)$, computed
    exactly as in the standard model.
    \item The \textbf{forward mapping} $\mathbf{g}(\btheta_i) = (\CE_1, \CE_2,
    \CE_3)$, a deterministic evaluation of the proposed parameters.
    \item The \textbf{normal prior density} $p_{\mathbf{Y}}(\mathbf{g}(\btheta_i)
    \given \bmu^Y, \bm{\sigma}^{Y2})$, evaluated at the CEs from step~2.
    \item The \textbf{Jacobian correction}
    $\log|\det J_{\mathbf{g}}(\btheta_i)|$, adjusting for the nonlinear
    change of variables.
\end{enumerate}
The log-posterior is:
\begin{equation}
    \label{eq:log_posterior_theta}
    \log p(\btheta_i \given \bmu^Y, \bm{\sigma}^{Y2}, \text{data}_i) =
    \log p(\text{data}_i \given \btheta_i) +
    \log p_{\mathbf{Y}}(\mathbf{g}(\btheta_i) \given \bmu^Y, \bm{\sigma}^{Y2})
    + \log \left|\det J_{\mathbf{g}}(\btheta_i)\right| + \text{const.}
\end{equation}
Crucially, \textbf{no inversion is required}. The forward mapping $\mathbf{g}$
and its Jacobian are available in closed form
(Section~\ref{sec:jacobian_rdu}). In practice, automatic differentiation
(e.g.\ \texttt{jax.jacfwd}) computes the $3 \times 3$ Jacobian matrix and its
determinant exactly, with negligible overhead relative to the likelihood
evaluation.

\begin{remark}[Why not sample in CE space?]
An alternative would draw proposals in $\mathbf{Y}$-space, where the prior is
normal. However, evaluating the likelihood then requires \emph{inverting} the
CE mapping---solving $\mathbf{g}(\btheta) = \mathbf{Y}$ at every MCMC step for
every subject---which is a $p$-dimensional root-finding problem. Since the
likelihood is evaluated thousands of times during sampling, this overhead is
substantial. Parameter-space sampling avoids inversion entirely: the forward
map and its Jacobian are all that is needed.
\end{remark}

\subsection{Why This Works}

With $p$ CEs for $p$ preference parameters and an invertible mapping, the
framework inherits all the simplicity of the identity case:
\begin{enumerate}
    \item \textbf{Correct dimensionality.} We have $2p$ moment conditions for
    $2p$ unknowns (after fixing $\rho$ and $\alpha_0$).

    \item \textbf{No intractable conditional moments.} The prior is specified
    directly in CE space, so the moments are just normal moments---no
    integration against a nonlinear $g$ is needed.

    \item \textbf{No coupling between mean and variance.} In CE space, the
    mean hyperparameters $\mu_k^Y$ and variance hyperparameters $\sigma_k^{Y2}$
    are separated exactly as in the identity case.

    \item \textbf{No approximation.} The Jacobian handles the nonlinearity of
    $\mathbf{g}$ exactly.
\end{enumerate}

\subsection{Requirements}

For this approach to work, three conditions must hold:

\paragraph{Dimensionality.} The number of calibration quantities must equal the
number of preference parameters: $\dim(\mathbf{Y}) = \dim(\btheta)$. Parameters
not involved in the observable mapping (such as the noise parameter $\lambda$ in
the RDU model) must be handled separately---either specified directly in
parameter space, or mapped through their own invertible function (e.g., a choice
probability for a specific lottery pair).

\paragraph{Invertibility.} The mapping $\mathbf{g}$ must be a bijection on the
relevant domain. This requires two ingredients: (i) \emph{local invertibility}
($\det J_{\mathbf{g}} \neq 0$ everywhere on the prior support), and
(ii) \emph{global injectivity} (no two distinct parameter vectors produce
identical CEs). For the RDU-Prelec CE mapping, local invertibility holds if
and only if the reference lotteries have distinct cumulative probability
structures (Section~\ref{sec:jacobian_rdu}). Global injectivity additionally
requires prize diversity across lotteries
(Section~\ref{sec:general_injectivity}); without it, utility curvature and
probability weighting can compensate each other. The general conditions for
global injectivity are developed in Section~\ref{sec:sufficient}.

\paragraph{Computability.} The forward mapping $\mathbf{g}$ and its Jacobian
$J_{\mathbf{g}}$ must be computable. For parameter-space sampling
(Section~\ref{sec:posterior_sampling}), the inverse $\mathbf{g}^{-1}$ is
\emph{not} required---only the forward evaluation
$\btheta \mapsto \mathbf{g}(\btheta)$ and the determinant
$|\det J_{\mathbf{g}}(\btheta)|$. Both are available analytically for the CE
mapping (Section~\ref{sec:jacobian_rdu}), and automatic differentiation
computes the Jacobian with negligible overhead.

\subsection{Application to RDU}
\label{sec:rdu_invertible}

For the RDU model with CRRA utility and 2-parameter Prelec PWF:

\paragraph{Preference parameters:} $\btheta = (r, \phi, \eta)$. These
determine the certainty equivalent of any lottery.

\paragraph{Noise parameter:} $\lambda$ governs the Fechner error structure and
does not affect the deterministic CE. It can be handled in parameter space
using the standard moment-matching approach, or mapped through a choice
probability for a specific lottery pair.

\paragraph{Reference lotteries:} Choose 3 lotteries $L_1, L_2, L_3$ with
diverse characteristics. For example:
\begin{itemize}
    \item $L_1$: Low stakes, equal probabilities (e.g., \$10/\$30/\$50,
        each $1/3$)---primarily sensitive to $r$.
    \item $L_2$: High stakes, equal probabilities (e.g., \$100/\$300/\$500,
        each $1/3$)---sensitive to $r$ at a different scale, allowing the
        nonlinearity of CRRA to separate $r$ from the PWF.
    \item $L_3$: Skewed probabilities (e.g., \$10 w.p.\ $0.9$, \$50 w.p.\
        $0.1$)---primarily sensitive to probability weighting ($\phi$, $\eta$).
\end{itemize}
The mapping is $\mathbf{g}(r, \phi, \eta) = (\CE_1, \CE_2, \CE_3)$.

\paragraph{Prior specification:} The researcher specifies, for each lottery:
\begin{itemize}
    \item $m^{k*}$: Expected CE for a random individual.
    \item $s^{2k*}$: Variance of CEs across individuals.
\end{itemize}
For a lottery paying \$10/\$30/\$50 equally, the expected value is \$30. A
researcher might specify $m^{1*} = 25$ (slight risk aversion) with
$s^{1*} = 8$ (standard deviation of \$8, so a 95\% interval of roughly
\$9--\$41). These are statements about dollars, not about abstract curvature
parameters.

\paragraph{Hyperprior derivation:} Apply the standard formulas in CE space:
\begin{align}
    \mu_{0,k}^Y &= m^{k*} \\
    \tau_{0,k}^{Y2} &= \rho \, s^{2k*} \\
    \beta_{0,k}^Y &= (1 - \rho) \, s^{2k*} \, (\alpha_0 - 1)
\end{align}
Identical to the current paper's equations, with CE moments in place of
parameter moments.

\paragraph{Correlation structure.} With independent priors across the 3 CEs,
the induced prior on $(r, \phi, \eta)$ will generally \emph{not} have
independent components---the nonlinearity of $\mathbf{g}^{-1}$ introduces
dependence. This is a feature: the CE-space specification implicitly encodes
cross-parameter structure that would be difficult to specify directly. If
desired, the researcher could specify a joint prior over CEs with correlations,
using the Normal-Inverse-Wishart framework from Appendix~A of the paper.

\subsection{Jacobian Structure and Invertibility Verification}
\label{sec:jacobian_rdu}

We derive the explicit Jacobian of the CE mapping for the RDU-Prelec model and
verify invertibility for a specific lottery set. The results in this section
were confirmed by symbolic computation (SymPy).

\paragraph{Monotone reduction.} Since
$\CE_k = [B_k]^{1/(1-r)}$ where $B_k = (1-r) \cdot \RDU_k$ is a monotone
transform of $\CE_k$ for fixed $r$, the mapping
$(r, \phi, \eta) \mapsto (\CE_1, \CE_2, \CE_3)$ is invertible if and only if
$(r, \phi, \eta) \mapsto (B_1, B_2, B_3)$ is. We work with $B_k$ for
algebraic simplicity.

For a 3-outcome lottery with prizes $x_1 < x_2 < x_3$, cumulative
probabilities $q_k = P(X \geq x_3)$ and $s_k = P(X \geq x_2)$, and decision
weights $d_1 = 1 - w(s_k)$, $d_2 = w(s_k) - w(q_k)$, $d_3 = w(q_k)$:
\begin{equation}
    \label{eq:B_k}
    B_k = (1 - w(s_k)) \, x_1^{1-r} + (w(s_k) - w(q_k)) \, x_2^{1-r}
    + w(q_k) \, x_3^{1-r}
\end{equation}

\paragraph{Jacobian entries.} Define the prize-power gaps
$\Delta_{12} \coloneqq x_2^{1-r} - x_1^{1-r} > 0$ and
$\Delta_{23} \coloneqq x_3^{1-r} - x_2^{1-r} > 0$. The partial derivatives of
$B_k$ are:
\begin{align}
    \frac{\partial B_k}{\partial r} &= -\Big[
        (1 - w(s_k)) \, x_1^{1-r} \ln x_1
        + (w(s_k) - w(q_k)) \, x_2^{1-r} \ln x_2
        + w(q_k) \, x_3^{1-r} \ln x_3 \Big]
    \label{eq:dBdr} \\[6pt]
    \frac{\partial B_k}{\partial \phi} &=
        \frac{\partial w}{\partial \phi}\bigg|_{s_k} \Delta_{12}
        + \frac{\partial w}{\partial \phi}\bigg|_{q_k} \Delta_{23}
    \label{eq:dBdphi} \\[6pt]
    \frac{\partial B_k}{\partial \eta} &=
        \frac{\partial w}{\partial \eta}\bigg|_{s_k} \Delta_{12}
        + \frac{\partial w}{\partial \eta}\bigg|_{q_k} \Delta_{23}
    \label{eq:dBdeta}
\end{align}
where the Prelec derivatives are:
\begin{align}
    \frac{\partial w}{\partial \phi}\bigg|_p &= -\eta \, (-\ln p)^\phi
        \ln(-\ln p) \, w(p)
    \label{eq:dwdphi} \\
    \frac{\partial w}{\partial \eta}\bigg|_p &= -(-\ln p)^\phi \, w(p)
    \label{eq:dwdeta}
\end{align}
Note that \eqref{eq:dBdphi} and \eqref{eq:dBdeta} have the same structure:
each is a linear combination of PWF derivatives evaluated at $s_k$ and $q_k$,
weighted by the prize-power gaps. The $r$-column \eqref{eq:dBdr} has a
fundamentally different structure involving $\ln x_j$ terms.

\paragraph{Prelec derivative ratio.} The ratio of the two Prelec derivatives is:
\begin{equation}
    \label{eq:prelec_ratio}
    \frac{\partial w / \partial \phi}{\partial w / \partial \eta}\bigg|_p
    = \eta \, \ln(-\ln p)
\end{equation}
This ratio \textbf{depends on $p$}. Substituting into
\eqref{eq:dBdphi}--\eqref{eq:dBdeta}, the ratio of the $\phi$-entry to the
$\eta$-entry in Jacobian row $k$ is a weighted average:
\begin{equation}
    \label{eq:column_ratio}
    \frac{\partial B_k / \partial \phi}{\partial B_k / \partial \eta}
    = \eta \Big[ \omega_k \ln(-\ln s_k) + (1 - \omega_k) \ln(-\ln q_k) \Big]
\end{equation}
where $\omega_k \in (0,1)$ depends on the prize gaps and PWF derivatives at
$s_k$. This weighted average varies across rows if and only if the lotteries
have different $(q_k, s_k)$ pairs.

\paragraph{Invertibility condition.} The Jacobian
$\partial(B_1, B_2, B_3) / \partial(r, \phi, \eta)$ is singular if and only if
the $\phi$ and $\eta$ columns are proportional, since the $r$ column is
structurally independent (involving $\ln x_j$ terms that the PWF columns lack).
By \eqref{eq:column_ratio}, proportionality holds if and only if all three
ratios \eqref{eq:column_ratio} are equal, which occurs precisely when
$q_1 = q_2 = q_3$ and $s_1 = s_2 = s_3$. Therefore:
\begin{quote}
    \emph{The CE mapping is generically invertible if and only if the reference
    lotteries have distinct cumulative probability structures $(q_k, s_k)$.}
\end{quote}
With identical probability structures, the $\phi$ and $\eta$ columns are
proportional with ratio $\eta \ln(-\ln q)$ (a scalar independent of $k$), and
$\det J = 0$ identically. This is confirmed by symbolic computation:
substituting equal probabilities yields $\det J = 0$ exactly.

\paragraph{Numerical verification of local invertibility.} For the no-FOSD
lottery set
\begin{center}
\begin{tabular}{lccc}
    Lottery & Probabilities $(x_1, x_2, x_3)$ & $q_k$ & $s_k$ \\
    \hline
    $L_1$ (binary)  & $(1/2,\; 0,\; 1/2)$   & $1/2$  & $1/2$ \\
    $L_2$ (middle)  & $(1/10,\; 8/10,\; 1/10)$ & $1/10$ & $9/10$ \\
    $L_3$ (skewed)  & $(3/10,\; 3/10,\; 4/10)$ & $4/10$ & $7/10$ \\
\end{tabular}
\end{center}
with shared prizes $(10, 30, 50)$, the CDFs cross for every pair (no pairwise
FOSD), guaranteeing that CE orderings are unconstrained across preference
types. Symbolic computation of $\det J$ evaluated at 8 points spanning
$r \in [0.1, 0.9]$, $\phi \in [0.4, 1.5]$, $\eta \in [0.4, 1.5]$ yields a
\emph{strictly positive} determinant at every point, confirming local
invertibility throughout the economically relevant domain.

\paragraph{Global injectivity fails for shared-prize lotteries.}
Despite $\det J > 0$ everywhere, the 3-lottery CE mapping with shared prizes
is \textbf{not globally injective}. Numerical computation identifies
observationally equivalent parameter triples---distinct
$(r, \phi, \eta) \neq (r', \phi', \eta')$ producing identical
$(\CE_1, \CE_2, \CE_3)$ to machine precision:
\begin{center}
\begin{tabular}{lcccc}
    & $r$ & $\phi$ & $\eta$ & max $|\Delta\CE|$ \\ \hline
    Params A & $0.80$ & $0.50$ & $1.30$ & \\
    Params B & $0.63$ & $0.52$ & $1.43$ & $< 10^{-9}$ \\[4pt]
    Params A & $1.50$ & $0.60$ & $1.00$ & \\
    Params B & $0.55$ & $0.71$ & $1.68$ & $< 10^{-9}$ \\
\end{tabular}
\end{center}
The two parameter sets produce \emph{different} decision weights for each
lottery, but the differences are exactly compensated by the change in utility
curvature: higher $r$ (more concave utility) is offset by lower $\eta$ and
$\phi$ (less pessimistic probability weighting), and vice versa.

\paragraph{The compensation mechanism.} With shared prizes, all lotteries
share the same utility differences $\Delta_{12}(r) = x_2^{1-r} - x_1^{1-r}$
and $\Delta_{23}(r) = x_3^{1-r} - x_2^{1-r}$. After normalization, the system
becomes:
\begin{equation}
    \label{eq:shared_prize_system}
    C_k(r) = w(s_k;\, \phi, \eta) + R(r) \cdot w(q_k;\, \phi, \eta),
    \quad k = 1, 2, 3
\end{equation}
where $R(r) = \Delta_{23}(r)/\Delta_{12}(r)$ is a \emph{single} scalar
function of $r$ common to all three equations. Although $R(r)$ is strictly
decreasing (proved: higher risk aversion compresses upper utility gaps more
than lower), the system \eqref{eq:shared_prize_system} has three equations in
three effective unknowns---$(R, \phi, \eta)$---and is therefore exactly
determined, admitting multiple solutions in $(r, \phi, \eta)$.

In economic terms: all three lotteries test probability weighting at different
probability levels, but they all test utility curvature at the \emph{same}
prize levels. The mapping has enough information to locally identify all
parameters, but globally there exists a smooth curve in parameter space along
which the utility and probability weighting effects exactly compensate.

\paragraph{Practical computation.} While the symbolic Jacobian entries
\eqref{eq:dBdr}--\eqref{eq:dBdeta} are available in closed form, in practice
one need not implement them by hand. Automatic differentiation
(\texttt{jax.jacfwd} or \texttt{jax.jacrev}) applied to the forward CE
function computes the full $3 \times 3$ Jacobian matrix---and hence its
log-determinant---exactly and efficiently. The computational cost is comparable
to a single forward evaluation of $\mathbf{g}$, which is itself far cheaper
than the likelihood evaluation over 80 choice observations per subject.

\subsection{Achieving Global Injectivity}
\label{sec:global_injectivity}

The shared-prize compensation mechanism suggests two natural remedies: adding
a fourth lottery (overdetermination) or using lotteries with distinct prize
structures (breaking the common factor).

\subsubsection{Strategy 1: A Fourth Lottery on the Same Prizes}

With $K = 4$ lotteries sharing prizes $(x_1, x_2, x_3)$, the system
\eqref{eq:shared_prize_system} becomes four equations in three effective
unknowns $(R, \phi, \eta)$---generically overdetermined with no solution for
the spurious root. Numerical verification confirms: every fourth lottery tested
(seven designs spanning diverse probability structures) separates both
confirmed equivalent pairs, with $|\Delta\CE| \in [0.04, 3.2]$.

This strategy is guaranteed to work for generic probability structures, since
the fourth equation imposes an additional constraint that the spurious
parameter triple generically cannot satisfy.

\subsubsection{Strategy 2: Lotteries with Distinct Prize Structures}

With different prizes across lotteries, the normalized system becomes:
\begin{equation}
    \label{eq:mixed_prize_system}
    C_k(r) = w(s_k;\, \phi, \eta) \cdot f_k(r)
    + w(q_k;\, \phi, \eta) \cdot g_k(r), \quad k = 1, 2, 3
\end{equation}
where $f_k(r)$ and $g_k(r)$ are lottery-specific functions of $r$ (determined
by the prize vectors). Now $r$ enters through \emph{multiple independent
channels} rather than through the single scalar $R(r)$. Matching all three
equations simultaneously overdetermines $r$, generically precluding
compensation.

Numerical verification is decisive: for six tested mixed-prize
configurations---including designs where all prizes remain within the
$[\$10, \$50]$ range---\textbf{no observationally equivalent triples were
found} for either confirmed equivalent pair. Even modest variation suffices:
replacing one lottery's prizes from $(10, 30, 50)$ to $(15, 30, 45)$ or
$(10, 20, 40)$ eliminates the equivalence entirely.

\paragraph{Which designs work.} The key requirement is that at least two
lotteries differ in their \emph{interior} prize spacing. The effective
condition is:
\begin{quote}
    \emph{At least two lotteries must have distinct ratios
    $\Delta_{23}^{(k)}(r) / \Delta_{12}^{(k)}(r)$ as functions of $r$.}
\end{quote}
Since $R_k(r) = (c_k^{1-r} - b_k^{1-r}) / (b_k^{1-r} - a_k^{1-r})$ depends
on the prize vector $(a_k, b_k, c_k)$, two lotteries with genuinely different
prize spacing produce different $R_k$ functions, and the compensation mechanism
breaks.

Importantly, \textbf{changing a single prize suffices}---the extreme prizes can
remain shared. For example, changing only the middle prize from \$30 to \$25
(i.e., using $(10, 25, 50)$ instead of $(10, 30, 50)$ for one lottery) is
enough to break the equivalence. The separation is modest but decisive:
$|\Delta\CE| \in [0.07, 0.78]$ across the confirmed equivalent pairs, and a
systematic search over $r \in [0.05, 3.0]$ with multiple solver restarts finds
\emph{no new equivalents} for any of four tested parameter configurations.

The only caveat is that the lottery must actually \emph{use} the changed
prize: a binary gamble placing zero probability on the middle outcome has its
CE determined entirely by the extreme prizes, so changing $b$ alone has no
effect for that specific lottery. As long as the middle outcome receives
positive probability, any change to $b$ produces a different $R_k(r)$ and
breaks the compensation.

\paragraph{Recommendation.} For the predictive prior specification, the
simplest design uses 3 lotteries where at least one differs in its middle
prize. For instance, within the existing $[\$10, \$50]$ experimental range:
\begin{itemize}
    \item $L_1$: Prizes $(10, 30, 50)$, diverse probabilities
    \item $L_2$: Prizes $(10, 30, 50)$, different probabilities
    \item $L_3$: Prizes $(10, 25, 50)$, any probabilities with $p_2 > 0$
\end{itemize}
The first two lotteries, sharing prizes, separate $\phi$ from $\eta$ through
different probability structures. The third lottery, differing only in its
middle prize, breaks the $r$-compensation through a different $R_3(r)$
function. Larger prize shifts (e.g., $(15, 30, 45)$ or $(10, 20, 40)$)
provide stronger separation but are not necessary.

\subsection{Relationship to the Current Framework}

The parameter-space approach in the current paper is the special case where the
reference quantities \emph{are} the parameters themselves:
$\mathbf{g} = \text{id}$, $J_{\mathbf{g}} = I$, and the Jacobian correction
is unity. The CE-space approach generalizes this by allowing the researcher to
choose any convenient coordinate system for specifying beliefs, with the
Jacobian mediating between the specification space and the likelihood space.

Importantly, $\rho$ and $\alpha_0$ retain their interpretations: $\rho$ is the
fraction of prior predictive CE variance attributable to uncertainty about the
population mean CE (versus heterogeneity in CEs across individuals), and
$\alpha_0$ controls the concentration of the CE variance distribution. These
are arguably \emph{more} interpretable in CE space than in parameter space,
since the researcher can reason about variance in dollars rather than variance
in utility curvature.


%==============================================================================
\section{Bounded Distributions for Observable-Space Priors}
\label{sec:bounded}
%==============================================================================

The Normal population distribution in Section~\ref{sec:invertible} is
convenient but incorrect for certainty equivalents. CEs are bounded:
\begin{equation}
    a_k \leq \CE_k \leq b_k
\end{equation}
where $a_k$ and $b_k$ are the minimum and maximum prizes of reference lottery
$k$. No agent values a lottery outside its prize range. This section develops
the proper bounded treatment using Beta distributions, shows that the variance
decomposition framework applies unchanged, and derives closed-form hyperprior
parameters analogous to the NIG case.

\subsection{The Full Hierarchy in Observable Space}

Before developing the bounded formulation, we state the hierarchy explicitly.
The key point is that \textbf{model parameters $\btheta$ never appear in the
prior specification}. They enter only through the likelihood and the Jacobian
that links observable and parameter space.

\begin{enumerate}
    \item[\textbf{Level 1.}] \textbf{Hyperprior.} A distribution over
    population-level parameters $(\mu_k, \phi_k)_{k=1}^p$ that govern the
    population distribution of each CE. These are denominated in observable
    units (dollars, normalized CEs) --- not in parameter space.

    \item[\textbf{Level 2.}] \textbf{Population distribution.} For subject
    $i$, the CE values are drawn from the population distribution:
    $z_i^k \given \mu_k, \phi_k \sim F(\mu_k, \phi_k)$, where $F$ is
    a distribution on $[0,1]$ (Beta in this section; Normal in the
    preceding section as a convenient approximation).

    \item[\textbf{Level 3.}] \textbf{Individual parameters.} The
    parameter vector $\btheta_i$ is related to the CEs through
    $\mathbf{g}(\btheta_i) = (Y_i^1, \ldots, Y_i^p)$. The prior density on
    $\btheta_i$ is the CE-space density evaluated at $\mathbf{g}(\btheta_i)$,
    multiplied by the Jacobian correction
    $\sqrt{\det(J_{\mathbf{g}}^\top J_{\mathbf{g}})}$.

    \item[\textbf{Level 4.}] \textbf{Data.} The likelihood
    $L(\text{data}_i \given \btheta_i)$ is evaluated in parameter space, as
    usual.
\end{enumerate}

\noindent
The variance decomposition operates at the interface between Levels 1 and 2:
the researcher specifies beliefs about unconditional CE moments (integrating
over the hyperprior), and $\rho$ and $\kappa$ partition the CE variance into
between-population and within-population components. The Jacobian correction
at Level 3 is a separate, independent step.

\subsection{Normalized Certainty Equivalents}

Define the normalized CE for reference lottery $k$:
\begin{equation}
    z_i^k = \frac{Y_i^k - a_k}{b_k - a_k} \in [0, 1]
\end{equation}
where $Y_i^k = g_k(\btheta_i)$ is the CE in dollar terms. The researcher
specifies beliefs in dollars and we convert to the normalized scale. For the
Harrison \& Ng (2016) data with prizes $\{10, 30, 50\}$: $z = 0$ corresponds
to $\CE = \$10$ (extreme risk aversion), $z = 0.5$ to $\CE = \$30$ (risk
neutrality), and $z = 1$ to $\CE = \$50$ (extreme risk seeking).

\subsection{Beta Population Distribution}
\label{sec:beta_pop}

The Beta distribution with mean-concentration parameterization is the natural
choice for $z \in [0,1]$:
\begin{equation}
    \label{eq:beta_pop}
    z_i^k \given \mu_k, \phi_k \sim
    \mathrm{Beta}(\mu_k \phi_k, \; (1 - \mu_k)\phi_k)
\end{equation}
where:
\begin{itemize}
    \item $\mu_k \in (0,1)$ is the population mean (in normalized units):
    the ``typical'' CE for lottery $k$, rescaled to $[0,1]$.
    \item $\phi_k > 0$ is the concentration parameter: higher $\phi_k$ means
    less within-population dispersion.
\end{itemize}
The moments are:
\begin{align}
    \E[z_i^k \given \mu_k, \phi_k] &= \mu_k \\
    \Var[z_i^k \given \mu_k, \phi_k] &= \frac{\mu_k(1-\mu_k)}{\phi_k + 1}
    \label{eq:beta_var}
\end{align}

Define the \emph{variance fraction}:
\begin{equation}
    \label{eq:psi_def}
    \psi_k \coloneqq \frac{1}{\phi_k + 1} \in (0,1)
\end{equation}
so that the within-population variance is
$\Var[z_i^k \given \mu_k, \phi_k] = \mu_k(1-\mu_k) \, \psi_k$. The quantity
$\psi_k$ has a direct interpretation: the within-population CE variance as a
fraction of its maximum possible value $\mu_k(1-\mu_k)$. When $\psi_k = 0$
(equivalently $\phi_k \to \infty$), the population is degenerate; when
$\psi_k = 1$ (equivalently $\phi_k = 0$), the population has maximum
dispersion (Bernoulli on the endpoints).

\subsection{Variance Decomposition}

The law of total variance applies unchanged. The unconditional variance of
$z_i^k$ decomposes as:
\begin{equation}
    \label{eq:beta_total_var}
    v_k^* = \Var[z_i^k]
    = \underbrace{\E\!\left[\frac{\mu_k(1-\mu_k)}{\phi_k + 1}\right]}_{\text{within-population heterogeneity}}
    + \underbrace{\Var[\mu_k]}_{\text{between-population (mean uncertainty)}}
\end{equation}

The researcher's two choices apply directly:
\begin{align}
    \Var[\mu_k] &= \rho \, v_k^*
    && \text{(between-population share)}
    \label{eq:beta_rho_between} \\
    \E\!\left[\frac{\mu_k(1-\mu_k)}{\phi_k + 1}\right] &= (1-\rho) \, v_k^*
    && \text{(within-population share)}
    \label{eq:beta_rho_within}
\end{align}

\begin{remark}[Distribution-free character of the decomposition]
Equations~\eqref{eq:beta_total_var}--\eqref{eq:beta_rho_within} have the same
structure as in the Normal-IG case. The law of total variance is a property of
iterated expectations, not of any particular distributional family. The $\rho$
decomposition is defined by the decomposition of unconditional CE variance, and
this decomposition is valid for any population distribution and any hyperprior.
What changes across families is the mapping from moment conditions to
hyperprior parameters.
\end{remark}

\subsection{Hyperprior Derivation}

\subsubsection{Hyperprior on $\mu_k$: the population mean CE}

Since $\mu_k \in (0,1)$, the natural hyperprior is itself a Beta:
\begin{equation}
    \mu_k \sim \mathrm{Beta}(a_\mu, b_\mu)
\end{equation}
The two moment conditions are:
\begin{align}
    \E[\mu_k] &= \mu_k^*
    && \Longrightarrow \quad \frac{a_\mu}{a_\mu + b_\mu} = \mu_k^*
    \\[4pt]
    \Var[\mu_k] &= \rho \, v_k^*
    && \Longrightarrow \quad
    \frac{\mu_k^*(1-\mu_k^*)}{a_\mu + b_\mu + 1} = \rho \, v_k^*
\end{align}
Define the ``effective sample size'' of the hyperprior on $\mu_k$:
\begin{equation}
    \label{eq:nu_mu}
    \nu_\mu \coloneqq a_\mu + b_\mu
    = \frac{\mu_k^*(1-\mu_k^*)}{\rho \, v_k^*} - 1
\end{equation}
Then:
\begin{equation}
    \label{eq:beta_hyper_mu}
    a_\mu = \mu_k^* \, \nu_\mu, \qquad
    b_\mu = (1-\mu_k^*) \, \nu_\mu
\end{equation}

\begin{remark}
The solution requires $\nu_\mu > 0$, i.e.,
$\rho \, v_k^* < \mu_k^*(1-\mu_k^*)$. This is always satisfied: the total
variance $v_k^*$ of any distribution on $[0,1]$ with mean $\mu_k^*$ cannot
exceed $\mu_k^*(1-\mu_k^*)$ (the Bernoulli bound), so
$\rho \, v_k^* < v_k^* \leq \mu_k^*(1-\mu_k^*)$ whenever $\rho < 1$.
\end{remark}

\subsubsection{Hyperprior on $\psi_k$: the variance fraction}

Using the variance fraction $\psi_k = 1/(\phi_k + 1) \in (0,1)$ from
\eqref{eq:psi_def}, the within-population variance is
$\mu_k(1-\mu_k) \cdot \psi_k$. Assuming $\mu_k \perp \psi_k$ (the population
mean and dispersion are a priori independent), the within-population target
\eqref{eq:beta_rho_within} becomes:
\begin{equation}
    (1-\rho) \, v_k^*
    = \E[\mu_k(1-\mu_k)] \cdot \E[\psi_k]
\end{equation}
Since $\E[\mu_k(1-\mu_k)] = \mu_k^*(1-\mu_k^*) - \Var[\mu_k]
= \mu_k^*(1-\mu_k^*) - \rho \, v_k^*$, solving for $\E[\psi_k]$:
\begin{equation}
    \label{eq:psi_target}
    \psi_k^* \coloneqq \E[\psi_k]
    = \frac{(1-\rho) \, v_k^*}{\mu_k^*(1-\mu_k^*) - \rho \, v_k^*}
\end{equation}

The second researcher choice---$\kappa$, the coefficient of variation of the
within-population variance---determines the second moment. Under the
approximation $\mu_k \approx \mu_k^*$ (valid when $\Var[\mu_k]$ is small
relative to $\mu_k^*(1-\mu_k^*)$):
\begin{equation}
    \kappa \approx \frac{\mathrm{SD}[\psi_k]}{\E[\psi_k]}
\end{equation}
This is the same interpretable quantity as in the Normal-IG case: $\kappa$
measures the researcher's confidence in their specification of
within-population heterogeneity. Small $\kappa$ means the researcher is
confident about the level of heterogeneity; large $\kappa$ means they are
uncertain.

A Beta hyperprior on $\psi_k$ is natural since $\psi_k \in (0,1)$:
\begin{equation}
    \psi_k \sim \mathrm{Beta}(a_\psi, b_\psi)
\end{equation}
Define the effective sample size:
\begin{equation}
    \label{eq:nu_psi}
    \nu_\psi \coloneqq a_\psi + b_\psi
    = \frac{1 - \psi_k^*}{\psi_k^* \, \kappa^2} - 1
\end{equation}
Then:
\begin{equation}
    \label{eq:beta_hyper_psi}
    a_\psi = \psi_k^* \, \nu_\psi, \qquad
    b_\psi = (1-\psi_k^*) \, \nu_\psi
\end{equation}

\begin{remark}[Comparison with the $\kappa$--$\alpha_0$ relationship]
In the Normal-IG case, $\kappa = 1/\sqrt{\alpha_0 - 2}$ relates $\kappa$ to
the Inverse-Gamma shape parameter. Here, $\kappa$ determines $\nu_\psi$, the
effective sample size of the Beta hyperprior on $\psi_k$. The functional form
differs, but the role is identical: $\kappa$ controls how concentrated the
hyperprior is around the target heterogeneity level. As $\kappa \to 0$,
$\nu_\psi \to \infty$ and $\psi_k$ concentrates at $\psi_k^*$; as
$\kappa \to \infty$, $\nu_\psi \to 0$ and the hyperprior becomes diffuse.
\end{remark}

\subsection{Summary: The Beta-Beta Hierarchy}

Collecting the results, the complete hierarchy for each reference lottery $k$
(in normalized CE units) is:

\begin{center}
\begin{tabular}{lll}
\toprule
\textbf{Level} & \textbf{Distribution} & \textbf{Parameters} \\
\midrule
Population &
    $z_i^k \given \mu_k, \psi_k \sim
    \mathrm{Beta}\!\left(\frac{\mu_k(1-\psi_k)}{\psi_k},\;
    \frac{(1-\mu_k)(1-\psi_k)}{\psi_k}\right)$ &
    $\mu_k \in (0,1)$, $\psi_k \in (0,1)$ \\[8pt]
Hyper (mean) &
    $\mu_k \sim \mathrm{Beta}(a_\mu, b_\mu)$ &
    $\nu_\mu = \mu_k^*(1-\mu_k^*)/(\rho v_k^*) - 1$ \\[4pt]
Hyper (dispersion) &
    $\psi_k \sim \mathrm{Beta}(a_\psi, b_\psi)$ &
    $\nu_\psi = (1-\psi_k^*)/(\psi_k^* \kappa^2) - 1$ \\
\bottomrule
\end{tabular}
\end{center}

\noindent
The researcher provides three inputs per CE:
\begin{itemize}
    \item $\mu_k^*$: prior mean CE (in normalized units)
    \item $v_k^*$: prior variance of CE (total, unconditional)
    \item Two shared choices: $\rho$ (variance partition) and $\kappa$
    (heterogeneity confidence)
\end{itemize}
These four quantities ($\mu_k^*$, $v_k^*$, $\rho$, $\kappa$) determine the
four hyperprior parameters ($a_\mu$, $b_\mu$, $a_\psi$, $b_\psi$) through
equations~\eqref{eq:nu_mu}--\eqref{eq:beta_hyper_mu} and
\eqref{eq:nu_psi}--\eqref{eq:beta_hyper_psi}. The structure is identical to
the Normal-IG case: two moments per CE, two researcher choices shared across
CEs.

\subsection{Connection to the Normal-IG Case}

The Normal-IG hierarchy of Section~\ref{sec:invertible} is a limiting case.
When the Beta concentration is large ($\phi_k \gg 1$, equivalently
$\psi_k \ll 1$) and the mean is well inside the bounds ($\mu_k^*$ not near 0
or 1):
\begin{equation}
    \mathrm{Beta}(\mu\phi, (1-\mu)\phi) \approx
    \Normal\!\left(\mu, \; \frac{\mu(1-\mu)}{\phi}\right)
    \qquad \text{for } \phi \gg 1
\end{equation}
In this regime the bounds are irrelevant and the two formulations agree. The
Beta formulation is preferred when:
\begin{itemize}
    \item The CE prior mean is near the bounds (e.g., highly risk-averse or
    risk-seeking populations).
    \item The CE prior standard deviation is a substantial fraction of the
    prize range.
    \item The researcher wishes the prior to encode the known bounded support
    as a structural constraint.
\end{itemize}

\subsection{Implementation}

The Beta formulation changes the population distribution but not the MCMC
mechanics. The log-posterior for subject $i$'s parameters $\btheta_i$ is:
\begin{multline}
    \label{eq:log_posterior_beta}
    \log p(\btheta_i \given \bmu, \bm{\psi}, \text{data}_i) =
    \log L(\text{data}_i \given \btheta_i) \\
    + \sum_{k=1}^p \log \mathrm{Beta}\!\left(z_i^k \;\Big|\;
    \frac{\mu_k(1-\psi_k)}{\psi_k},\;
    \frac{(1-\mu_k)(1-\psi_k)}{\psi_k}\right)
    + \frac{1}{2}\log\det\!\big(J^\top J\big)
    + \text{const.}
\end{multline}
where $z_i^k = (g_k(\btheta_i) - a_k)/(b_k - a_k)$ is the normalized CE from
the forward map. No conjugacy is exploited---the MCMC sampler evaluates each
log-density term directly. The Jacobian correction is identical to the Normal
case; only the population-level density evaluation changes.

For the hyperparameter blocks, the Beta hyperpriors on $\mu_k$ and $\psi_k$
are evaluated directly:
\begin{equation}
    \log p(\mu_k, \psi_k) = \log \mathrm{Beta}(\mu_k \given a_\mu, b_\mu)
    + \log \mathrm{Beta}(\psi_k \given a_\psi, b_\psi)
\end{equation}
The cost per MCMC step is negligible: evaluating a Beta log-density is a single
call to \texttt{jax.scipy.stats.beta.logpdf}.


%==============================================================================
\section{Global Injectivity: General Conditions}
\label{sec:general_injectivity}
%==============================================================================

The shared-prize non-injectivity result for RDU-Prelec is an instance of a
general phenomenon. This section derives the conditions under which a CE
mapping is globally injective for an arbitrary utility model. The results
clarify what the researcher must verify for any specific application.

\subsection{Setup}

Consider a general decision-under-uncertainty model with parameter vector
$\bm{\alpha} = (\alpha_1, \ldots, \alpha_p) \in \Omega \subseteq \mathbb{R}^p$
and $K$ reference lotteries. Each lottery $k$ has outcomes
$\mathbf{x}_k = (x_{k,1}, \ldots, x_{k,n_k})$ with probabilities
$\mathbf{p}_k = (p_{k,1}, \ldots, p_{k,n_k})$. The certainty equivalent is:
\begin{equation}
    \label{eq:general_ce}
    \CE_k(\bm{\alpha}) = u^{-1}\!\left(
    \sum_{j=1}^{n_k} d_{k,j}(\bm{\alpha}) \, u(x_{k,j};\, \bm{\alpha})
    \,;\; \bm{\alpha}\right)
\end{equation}
where $u(x;\, \bm{\alpha})$ is the utility function and
$d_{k,j}(\bm{\alpha})$ are decision weights (which may depend on
$\bm{\alpha}$ through a probability weighting function, or may simply equal
$p_{k,j}$ for expected utility). Define the CE mapping
$\mathbf{g}: \Omega \to \mathbb{R}^K$ by
$\mathbf{g}(\bm{\alpha}) = (\CE_1(\bm{\alpha}), \ldots, \CE_K(\bm{\alpha}))$.

For the Jacobian-based prior specification of Section~\ref{sec:invertible}, we
require $K = p$ and $\mathbf{g}$ to be a diffeomorphism (smooth bijection).
Two conditions are needed: \emph{local} invertibility ($\det J \neq 0$
everywhere) and \emph{global} injectivity (no two distinct parameter vectors
produce identical CEs).

\subsection{Necessary Condition: Local Invertibility}

\begin{proposition}[Local invertibility]
\label{prop:local}
The CE mapping $\mathbf{g}: \Omega \to \mathbb{R}^p$ is a local
diffeomorphism if and only if $\det J_{\mathbf{g}}(\bm{\alpha}) \neq 0$ for
all $\bm{\alpha} \in \Omega$.
\end{proposition}

This is the inverse function theorem. It guarantees that in a neighborhood of
every point, the mapping can be inverted. For practical applications, one
verifies this either symbolically (as in Section~\ref{sec:jacobian_rdu}) or
numerically by evaluating $\det J$ on a grid and confirming it does not change
sign.

In 1 dimension ($p = 1$), local invertibility is \emph{sufficient} for global
injectivity: $g'(\alpha) > 0$ everywhere implies $g$ is strictly increasing.
In higher dimensions this fails, as the following makes precise.

\subsection{The Obstruction: Why Local Does Not Imply Global}
\label{sec:obstruction}

A local diffeomorphism $\mathbf{g}: \Omega \to \mathbb{R}^p$ (with $\Omega$
convex, $\det J > 0$ everywhere) can fail to be globally injective. The
canonical counterexample is $f(x,y) = (e^x \cos y,\; e^x \sin y)$ on
$\mathbb{R}^2$: $\det J = e^{2x} > 0$ everywhere, yet
$f(0, 0) = f(0, 2\pi)$.

For CE mappings, the obstruction has a specific economic interpretation.
The standard sufficient condition for promoting local to global invertibility
is \emph{properness}: preimages of compact sets must be compact (Hadamard's
global inverse function theorem). But CE mappings are never proper: the
codomain is bounded (CEs lie between the minimum and maximum prizes) while the
parameter domain is typically unbounded. As parameters approach extremes, CEs
saturate at boundary values, creating the possibility---in principle---that
distant parameter vectors map to identical CEs.

Whether this possibility is realized depends on the \emph{structure} of the
model, as the RDU-Prelec analysis demonstrates.

\subsection{The Compensation Mechanism}

The RDU-Prelec counterexample reveals a general obstruction to global
injectivity. Suppose the parameter vector decomposes as
$\bm{\alpha} = (\bm{\xi}, \bm{\psi})$ where $\bm{\xi}$ governs the utility
function $u(x;\, \bm{\xi})$ and $\bm{\psi}$ governs the decision weights
$d_{k,j}(\bm{\psi})$. Then:
\begin{equation}
    \CE_k(\bm{\xi}, \bm{\psi}) = u^{-1}\!\left(
    \sum_{j} d_{k,j}(\bm{\psi}) \, u(x_{k,j};\, \bm{\xi}) \,;\; \bm{\xi}\right)
\end{equation}
When all lotteries share the same prize vector $\mathbf{x}$, the utility
values $u(x_j;\, \bm{\xi})$ are \emph{common} across lotteries. The CE
mapping then has the form:
\begin{equation}
    \CE_k = u^{-1}\!\big(\mathbf{d}_k(\bm{\psi})^\top
    \mathbf{u}(\bm{\xi}) \,;\; \bm{\xi}\big)
\end{equation}
where $\mathbf{u}(\bm{\xi}) = (u(x_1;\,\bm{\xi}), \ldots, u(x_n;\,\bm{\xi}))$
is the same for all $k$. The parameter $\bm{\xi}$ enters only through the
shared ``utility profile'' $\mathbf{u}$, creating a common factor that
allows changes in $\bm{\xi}$ to be compensated by changes in $\bm{\psi}$.

With different prize structures, each lottery has its own utility profile
$\mathbf{u}_k(\bm{\xi})$, and $\bm{\xi}$ enters through multiple independent
channels, generically breaking the compensation.

\subsection{Sufficient Conditions for Global Injectivity}
\label{sec:sufficient}

\begin{proposition}[Overdetermination]
\label{prop:overdetermined}
If $K > p$ reference lotteries are used, each with
$\det J_{\mathbf{g}} \neq 0$ for any $p$-element subset, then the CE mapping
$\mathbf{g}: \Omega \to \mathbb{R}^K$ is generically globally injective
(the set of $\bm{\alpha}$ admitting an observationally equivalent
$\bm{\alpha}'$ has measure zero in $\Omega$).
\end{proposition}

This is the ``fourth lottery'' strategy: a single additional lottery
generically breaks any exact compensation between parameters.

\begin{proposition}[Prize diversity]
\label{prop:prize_diversity}
Consider a model with utility parameters $\bm{\xi} \in \mathbb{R}^{p_1}$ and
weighting parameters $\bm{\psi} \in \mathbb{R}^{p_2}$ ($p_1 + p_2 = p$). The
CE mapping with $K = p$ lotteries is globally injective (for generic
probability structures) if:
\begin{enumerate}
    \item[\textnormal{(a)}] The lotteries have distinct cumulative probability
    structures (separates the components of $\bm{\psi}$);
    \item[\textnormal{(b)}] At least $p_1 + 1$ distinct prize vectors appear
    among the lotteries (overdetermines $\bm{\xi}$ given $\bm{\psi}$);
    \item[\textnormal{(c)}] $\det J_{\mathbf{g}}(\bm{\alpha}) \neq 0$ for all
    $\bm{\alpha} \in \Omega$.
\end{enumerate}
\end{proposition}

Condition (b) is the key structural requirement. For the RDU-Prelec model with
$p_1 = 1$ (only $r$ is a utility parameter), at least $2$ distinct prize
vectors are needed. Since $K = 3$ lotteries are required, the condition allows
two lotteries to share prizes and the third to differ---exactly the design
verified numerically in Section~\ref{sec:global_injectivity}.

\paragraph{Intuition.} The argument parallels instrument relevance in
econometrics. With shared prizes, $\bm{\xi}$ enters through a single ``reduced
form'' (the utility profile $\mathbf{u}$), so the system is underidentified in
$\bm{\xi}$. Prize diversity introduces multiple ``instruments'' for
$\bm{\xi}$, each constraining it from a different direction.

\paragraph{Proof sketch for $p_1 = 1$.} Suppose lotteries $k$ and $j$ have
different prize vectors. The ratio
$R_k(r) = \Delta_{23}^{(k)}(r) / \Delta_{12}^{(k)}(r)$ and
$R_j(r) = \Delta_{23}^{(j)}(r) / \Delta_{12}^{(j)}(r)$ are both strictly
monotone in $r$, but with different functional forms (determined by their prize
vectors). For fixed $(\phi, \eta)$, each lottery's CE equation
$C_k(r) = w(s_k) + R_k(r) \cdot w(q_k)$ determines $r$ as an implicit
function. Two equations with different $R_k(\cdot)$ generically determine $r$
uniquely, leaving no freedom for the ``second solution'' that the shared-prize
case admits.

\subsection{The Pure Utility Case: No Probability Weighting}

When the decision weights are known (i.e., $d_{k,j} = p_{k,j}$ as in expected
utility, or more generally any model where the weights do not depend on
$\bm{\alpha}$), the CE mapping simplifies to:
\begin{equation}
    \CE_k(\bm{\alpha}) = u^{-1}\!\bigg(\sum_{j} p_{k,j} \, u(x_{k,j};\,
    \bm{\alpha})\bigg)
\end{equation}
with $\bm{\alpha} = \bm{\xi}$ (all parameters are utility parameters, $p_2 = 0$).
This case is fundamentally easier because there is no compensation
mechanism---all parameters enter through the utility function alone.

\paragraph{Single parameter ($p = 1$).}
If $\bm{\alpha} = r$ is a scalar (e.g., CRRA), then $\CE_k(r)$ is a scalar
function of $r$. Provided $\partial\CE_k / \partial r \neq 0$ (the lottery is
non-degenerate and $r$ affects the CE), strict monotonicity gives global
injectivity from a single lottery. No prize diversity is needed.

\paragraph{Multiple utility parameters ($p > 1$).}
If $\bm{\alpha} = (\alpha_1, \ldots, \alpha_p)$ are all utility parameters
(e.g., a utility function with curvature, loss aversion, and reference point
parameters), the CE mapping $\mathbf{g}: \mathbb{R}^p \to \mathbb{R}^p$
requires $K = p$ reference lotteries. The conditions for global injectivity
are:
\begin{enumerate}
    \item \textbf{Local invertibility}: $\det J_{\mathbf{g}} \neq 0$
    everywhere. This requires the $p$ lotteries to have sufficiently diverse
    characteristics that each parameter has a distinct ``signature'' in the
    Jacobian---verified symbolically or numerically.

    \item \textbf{Prize diversity}: At least $p + 1$ distinct prize
    vectors\footnote{More precisely: the utility profiles
    $\mathbf{u}_k(\bm{\alpha}) = (u(x_{k,1};\,\bm{\alpha}), \ldots,
    u(x_{k,n_k};\,\bm{\alpha}))$ must span a space of dimension exceeding
    $\dim(\bm{\alpha})$ as functions of $\bm{\alpha}$. With shared prizes, all
    profiles are the \emph{same} vector $\mathbf{u}(\bm{\alpha})$, giving
    dimension 1 regardless of $p$.} among the $K$ lotteries. This prevents the
    shared-utility-profile obstruction: with shared prizes, all $p$ parameters
    enter through a single utility vector, and the system of CE equations has
    the structure of $p$ equations in what is effectively a lower-dimensional
    image.
\end{enumerate}

\noindent
For the pure utility case, the compensation mechanism of
Section~\ref{sec:obstruction} is between different utility parameters (e.g.,
curvature compensating loss aversion) rather than between utility and weighting
parameters. The remedy is the same: prize diversity ensures each parameter
influences the CEs through independent channels.

\subsection{Summary of Conditions}

For a CE mapping $\mathbf{g}: \bm{\alpha} \mapsto (\CE_1, \ldots, \CE_K)$
with $K = p = \dim(\bm{\alpha})$:
\begin{center}
\begin{tabular}{lll}
    \textbf{Condition} & \textbf{What it ensures} & \textbf{Status} \\ \hline
    $\det J \neq 0$ everywhere & Local invertibility & Necessary \\
    Distinct probability structures & Separates $\bm{\psi}$ components &
        Necessary (with PWF) \\
    Prize diversity (condition~(b)) & Breaks $\bm{\xi}$--$\bm{\psi}$
        compensation & Sufficient for global \\
    Overdetermination ($K > p$) & Breaks \emph{all} compensation &
        Sufficient for global \\
\end{tabular}
\end{center}
The first two conditions are necessary for local invertibility. The third or
fourth is needed to promote local to global injectivity. For the parameter-space
Jacobian approach (Section~\ref{sec:posterior_sampling}), local invertibility
alone suffices for a well-defined density---but the researcher should be aware
that the induced prior may assign positive probability to observationally
equivalent regions unless global injectivity holds.



%==============================================================================
\section{Detailed Analysis for the RDU Application}
%==============================================================================

\subsection{The Certainty Equivalent as \texorpdfstring{$g(\btheta)$}{g(theta)}}

For the RDU model with a 3-outcome lottery $(x_1, x_2, x_3)$ with equal
probabilities, the certainty equivalent is:
\begin{equation}
    \CE(r, \phi, \eta) = u^{-1}\!\Big(\sum_{k=1}^3 d_k(r, \phi, \eta) \,
    u(x_k; r)\Big)
\end{equation}
where $u(x; r) = x^{1-r}/(1-r)$ is the CRRA utility function, and
$d_k(\phi, \eta)$ are the Prelec decision weights. (We suppress the noise
parameter $\lambda$ here since the CE is a preference-theoretic quantity
independent of the choice process.)

This function has several properties relevant to our analysis:

\paragraph{Gradient structure.} The partial derivatives of $\CE$ with respect
to the parameters capture different aspects of preferences:
\begin{itemize}
    \item $\partial \CE / \partial r$: Sensitivity to risk aversion. Negative
    for gains (higher $r$ means lower CE). Large in magnitude when the lottery
    has high variance.
    \item $\partial \CE / \partial \phi$: Sensitivity to probability distortion
    curvature. Captures the inverse-S shape effect.
    \item $\partial \CE / \partial \eta$: Sensitivity to probability distortion
    elevation. Captures overall optimism/pessimism.
\end{itemize}
These gradients are functions of $\btheta$ itself, evaluated at the population
mean. They are heterogeneous across parameters, meaning different parameters
contribute differently to CE variance.

\paragraph{Nonlinearity.} The CRRA utility function has a pole at $r = 1$
(where it transitions to log utility) and is highly nonlinear for $r$ near 1.
The Prelec function $w(p) = \exp(-\eta(-\ln p)^\phi)$ is also nonlinear in
$\phi$ and $\eta$. These nonlinearities make the Jacobian correction
essential---a linear approximation would be unreliable for large parameter
variances.

\paragraph{Parameter scale mismatch.} In the current implementation, $\phi$ and
$\eta$ are sampled on the log scale ($\phi_{\text{raw}} = \ln \phi$). The
function $g$ is then a composition: $\CE(r, e^{\phi_{\text{raw}}},
e^{\eta_{\text{raw}}})$. The exponential transform amplifies the effect of
variance in the raw parameters, especially in the tails.



%==============================================================================
\section{Implications for the Paper}
%==============================================================================

\subsection{What Changes}

The invertible-$\mathbf{g}$ approach (Section~\ref{sec:invertible}) is not
merely an extension of the current framework---it is the \emph{same} framework
applied in a different coordinate system. The paper's methodology section could
be presented at two levels:

\paragraph{General statement.} The moment-matching framework applies to any
quantity with defined moments. Given a direct prior over $\mathbf{Y} =
\mathbf{g}(\btheta)$---whether $\mathbf{g}$ is the identity (parameters
directly) or an economically meaningful mapping (certainty equivalents)---the
law of total variance decomposes the prior predictive variance, and $\rho$
partitions this variance between population-mean uncertainty and
within-population heterogeneity. The hyperprior is derived by matching moments
in whichever coordinate system the researcher specified beliefs.

\paragraph{The identity case as baseline.} The current paper's derivation is
the special case $\mathbf{g} = \text{id}$, which requires no Jacobian and
admits particularly clean closed-form hyperparameters.

\paragraph{The predictive case as generalization.} When
$\mathbf{g} \neq \text{id}$ but is invertible, the same formulas apply in
$\mathbf{g}$-space. Sampling proceeds in parameter space with the prior density
evaluated via the forward mapping $\mathbf{g}(\btheta)$ and the Jacobian
correction \eqref{eq:jacobian_prior}. No inversion or approximation is needed.

\subsection{What Is Preserved}

\begin{itemize}
    \item The $\rho$ and $\alpha_0$ researcher choices retain their
    interpretations and apply identically in any coordinate system.

    \item The moment-matching derivation of hyperpriors is unchanged---it is
    applied in the coordinate system where beliefs are specified.

    \item The hyperprior families (NIG, NIW) and their closed-form solutions
    are untouched.

    \item The ``no free lunch'' result holds: the number of degrees of freedom
    is the same; they are just expressed in a different space.
\end{itemize}

\subsection{Recommended Framing for the Paper}

The invertible case suggests a stronger framing than the original version of
this note proposed:

\begin{enumerate}
    \item Present the general framework in coordinate-free terms: the
    researcher specifies a direct prior over any $p$-dimensional quantity
    $\mathbf{Y} = \mathbf{g}(\btheta)$, and the moment-matching machinery
    derives hyperpriors in $\mathbf{Y}$-space.

    \item Develop the parameter-space case ($\mathbf{g} = \text{id}$) as the
    primary worked example, since it requires no inversion and yields the
    cleanest formulas.

    \item Note that when $\mathbf{g}$ is an invertible mapping to observable
    quantities (such as certainty equivalents), the framework applies directly
    in that space. Sampling in parameter space requires only the forward
    mapping $\mathbf{g}(\btheta)$ and a Jacobian correction---no inversion.

    \item Frame the existing prior predictive CE analysis as evidence that the
    parameter-level specification implies reasonable predictive behavior---and
    note that the researcher could equally have started from predictive beliefs.
\end{enumerate}

This positions the paper's contribution as more general than
``hyperprior selection for Normal-Inverse-Gamma BHMs''---it is a coordinate-free
moment-matching framework that applies in any parameterization the researcher
finds natural.


%==============================================================================
\section{Connections to the Literature}
%==============================================================================

\subsection{Prior Predictive Checking}

The prior predictive distribution $p(Y) = \int p(Y \given \theta) p(\theta)
\, d\theta$ has been advocated as a tool for prior assessment by Box (1980),
Gelman et al.\ (2013, Ch.~6), and Gabry et al.\ (2019, ``Visualization in
Bayesian workflow''). The standard approach is:
\begin{enumerate}
    \item Specify a prior.
    \item Simulate from the prior predictive.
    \item Check whether the implied distribution of observables is reasonable.
    \item If not, revise the prior.
\end{enumerate}
The extension proposed here collapses steps 1 and 3: the prior is
\emph{defined} by the predictive targets, rather than checked against them
post-hoc. This is conceptually cleaner but computationally harder.

\subsection{Calibrated Priors}

The idea of choosing priors to match predictive implications has been explored
in several contexts:
\begin{itemize}
    \item \textbf{Penalized complexity priors} (Simpson et al., 2017):
    Priors are calibrated so that the model assigns a specified probability to
    deviating from a simpler base model. This is a form of predictive
    calibration using KL divergence rather than moment matching.

    \item \textbf{Prior elicitation via observables} (Kadane and Wolfson, 1998;
    O'Hagan et al., 2006): Experts specify credible intervals for observable
    quantities; priors are derived by inversion. This is closest in spirit to
    the extension proposed here, but the elicitation literature focuses on
    conjugate models where the inversion is tractable.

    \item \textbf{Simulation-based calibration} (Talts et al., 2018): Uses
    simulation to validate that the prior-likelihood-posterior pipeline is
    self-consistent. This is a diagnostic tool rather than a prior specification
    method, but the forward simulation machinery is the same.
\end{itemize}

\subsection{Decision-Theoretic Prior Selection}

In the risk preference context specifically, the certainty equivalent is the
natural ``decision-relevant'' quantity. Specifying priors via CE targets
connects to the decision-theoretic literature on loss functions and utility:
the prior is calibrated so that the \emph{decision-relevant implications} of
parameter uncertainty are reasonable, not just the parameter values themselves.

This is especially compelling for policy-relevant applications where the
ultimate quantity of interest is a welfare measure (willingness to pay,
certainty equivalent, compensating variation) rather than a structural
parameter.


%==============================================================================
\section{Open Questions}
%==============================================================================

\begin{enumerate}
    \item \textbf{Global injectivity and lottery design (resolved).}
    Section~\ref{sec:jacobian_rdu} establishes local invertibility ($\det J > 0$
    everywhere) for distinct probability structures.
    Section~\ref{sec:global_injectivity} demonstrates that local invertibility
    does \emph{not} imply global injectivity for lotteries sharing a common
    prize vector, due to the $\bm{\xi}$--$\bm{\psi}$ compensation mechanism.
    Section~\ref{sec:sufficient} provides sufficient conditions for global
    injectivity: either overdetermination ($K > p$) or prize diversity. The
    practical recommendation is clear: use at least two distinct prize vectors
    among the reference lotteries.

    \item \textbf{Sensitivity to lottery choice.} Different reference lotteries
    define different coordinate systems in CE space, and therefore different
    induced priors in parameter space (even from the same CE-space prior). How
    sensitive are the posterior results to this choice? In principle, the
    posterior should be invariant to the parameterization (the Jacobian ensures
    consistency), but in practice, different CE-space priors represent different
    substantive beliefs.

    \item \textbf{Normality in CE space (resolved).} The Normal population
    distribution places mass outside the CE bounds $[a_k, b_k]$.
    Section~\ref{sec:bounded} resolves this by replacing the Normal with a
    Beta distribution on the normalized CE $z \in [0,1]$, using the
    mean-concentration parameterization. The variance decomposition framework
    applies unchanged, and the Beta-Beta hyperprior has closed-form parameters
    determined by $(\mu_k^*, v_k^*, \rho, \kappa)$---the same inputs as the
    Normal-IG case.

    \item \textbf{Correlation structure.} Independent normal priors across CEs
    induce a specific (non-independent, non-normal) prior over $(r, \phi, \eta)$
    via $\mathbf{g}^{-1}$. Is this induced dependence structure reasonable?
    Conversely, does the researcher have intuition about CE correlations (e.g.,
    ``someone with a low CE for a small-stakes lottery also has a low CE for a
    large-stakes lottery'') that should be encoded?

    \item \textbf{Handling the noise parameter.} $\lambda$ does not affect CE
    but does affect choice probabilities. It could be: (a) specified directly in
    parameter space, (b) mapped through a choice probability (e.g., $P(\text{choose
    safer option})$ for a specific lottery pair), or (c) integrated out. Option
    (b) would bring the full model into predictive space but requires a
    4-dimensional invertible mapping.
\end{enumerate}



%==============================================================================
\section{Summary}
%==============================================================================

The moment-matching framework for hyperprior selection is
coordinate-free: the law of total variance and the $\rho$ decomposition apply
identically to any quantity with defined moments, whether that quantity is a
structural parameter or an observable like a certainty equivalent.

The key insight of this note is that when $\mathbf{g}: \mathbb{R}^p \to
\mathbb{R}^p$ is an invertible mapping from preference parameters to observable
quantities, the \emph{entire} Bayesian hierarchy---from hyperprior through
population distribution---can be specified in observable space with no
approximation. Model parameters never enter the prior specification; they
appear only in the likelihood and the Jacobian correction that links the two
spaces. The researcher specifies beliefs about certainty equivalents
(or other interpretable quantities), the standard $\rho$/$\kappa$
decomposition derives hyperpriors in that space, and the Jacobian of
$\mathbf{g}$ handles the coordinate change exactly. The identity case
($\mathbf{g} = \text{id}$, parameter-level specification) is a special case
with unit Jacobian.

Because certainty equivalents are bounded by the prize range, the
Normal population distribution should be replaced by a Beta distribution on
normalized CEs (Section~\ref{sec:bounded}). The Beta-Beta hierarchy---Beta
population distribution with Beta hyperpriors on the mean and variance
fraction---has closed-form hyperprior parameters determined by the same four
inputs as the Normal-IG case: the prior mean $\mu_k^*$, the prior variance
$v_k^*$, the variance partition $\rho$, and the heterogeneity confidence
$\kappa$. The variance decomposition is distribution-free; only the mapping
from moment conditions to hyperprior parameters changes across families.

A central finding of this note is that \textbf{local invertibility ($\det J
\neq 0$ everywhere) does not guarantee global injectivity} for the CE mapping.
When all reference lotteries share a common prize vector, utility curvature and
probability weighting can exactly compensate each other, producing distinct
parameter vectors with identical certainty equivalents. This compensation
mechanism is broken by either using lotteries with distinct prize structures or
by adding an additional reference lottery (overdetermination). The general
conditions for global injectivity---applicable to any utility model with
separable utility and weighting parameters---are developed in
Section~\ref{sec:general_injectivity}.

For the current paper, the invertible case strengthens the contribution by
showing that the framework is not tied to parameter-space specification. The
parameter-level derivation can be presented as the cleanest special case, with
the CE-space generalization noted as immediate. For a follow-up paper or an
expanded version, a full worked example in CE space---with reference lottery
selection, Jacobian computation, and posterior comparison to the parameter-space
approach---would demonstrate the practical value of the extension.


\end{document}
